\chapter{Inserting an Internal Cardioverter Defibrillator}

\section*{Definition and Overview}
An implantable cardioverter defibrillator (ICD) is a small device placed under the skin of the chest that continuously monitors the heart rhythm and delivers a shock to restore a normal rhythm if a life-threatening fast arrhythmia occurs. It is used in people who have survived a cardiac arrest or dangerous arrhythmia, or who are at high risk of one because of severe heart failure, certain inherited heart conditions, or poor heart pumping function. Insertion is similar to pacemaker implantation, performed under local anaesthetic with sedation, with leads passed through veins into the heart. The ICD does not cure the underlying heart condition but provides a safety net that can save life.

\section*{Epidemiology}
Around 15,000 to 20,000 Australians experience an out-of-hospital cardiac arrest each year, and most are due to ventricular arrhythmias. ICDs are implanted in many thousands of Australians, and the rate of implantation has risen as evidence shows they improve survival in people with heart failure and reduced pumping function. Most recipients have ischaemic heart disease or cardiomyopathy, and a smaller number have inherited conditions such as long QT syndrome, Brugada syndrome, or hypertrophic cardiomyopathy. As the population ages and heart failure becomes more common, ICD use continues to grow.

\section*{Aetiology and Pathophysiology}
Life-threatening arrhythmias, such as ventricular tachycardia and ventricular fibrillation, occur when the heart's electrical system becomes unstable, often on a background of scarred or diseased heart muscle from heart attacks, cardiomyopathy, or inherited ion channel disorders. During ventricular fibrillation the heart quivers instead of pumping, and without defibrillation within minutes, death follows. An ICD detects these rhythms and terminates them by delivering a shock, restoring coordinated pumping. Modern devices also act as pacemakers, and many are combined with cardiac resynchronisation therapy to treat heart failure. Implanting an ICD is an effective secondary prevention (after an event) or primary prevention (in high-risk people) strategy.

\section*{Clinical Features}
\begin{itemize}
    \item Indications: a previous cardiac arrest, documented dangerous arrhythmia, or severe heart failure with low pumping function
    \item Insertion under local anaesthetic with sedation, taking about an hour
    \item A small scar over the collarbone and restriction of arm movement for a few weeks
    \item The device is checked before discharge and at intervals afterwards
    \item If a dangerous rhythm occurs, the ICD delivers a shock, which the person may feel as a thump
    \item The device also records events and can be checked remotely
    \item Battery life is typically 5 to 10 years
\end{itemize}

\section*{Differential Diagnosis}
The decision to implant an ICD follows assessment of the person's arrhythmia risk and overall health.
\begin{itemize}
    \item \textbf{Heart failure management:} medicines and devices that improve function before considering an ICD
    \item \textbf{Reversible causes of arrhythmia:} heart attacks, electrolyte disturbances, and medicines
    \item \textbf{Coronary revascularisation:} improving blood supply to the heart
    \item \textbf{Inherited heart conditions:} assessment and family screening
    \item \textbf{Structural heart disease:} valve disease and cardiomyopathy
    \item \textbf{Life expectancy and quality of life:} whether the device aligns with the person's goals
    \item \textbf{Alternatives:} a wearable defibrillator or subcutaneous ICD in selected cases
\end{itemize}

\section*{When to Seek Medical Care}
People with heart disease should have their arrhythmia risk assessed, and those considered at high risk should discuss an ICD with their cardiologist. Anyone who has survived a cardiac arrest or a dangerous arrhythmia should be referred urgently. After implantation, people should attend device checks and seek review for any problem or shock.

\textbf{Call 000 (or local emergency services) for:}
\begin{itemize}
    \item Fainting, collapse, or a prolonged loss of consciousness
    \item A single shock followed by feeling unwell, or repeated shocks
    \item Chest pain, breathlessness, or palpitations
    \item Signs of infection: redness, swelling, pain, or fever at the device site
    \item Any cardiac arrest: start CPR and use a public defibrillator while waiting
\end{itemize}

\section*{Diagnosis and Investigations}
\begin{itemize}
    \item \textbf{Electrocardiogram and rhythm monitoring:} to detect and characterise arrhythmias
    \item \textbf{Echocardiogram:} assessment of pumping function and structure
    \item \textbf{Coronary angiography:} to assess for ischaemic heart disease
    \item \textbf{Cardiac MRI:} in selected cases for cardiomyopathy assessment
    \item \textbf{Electrophysiological studies:} testing of the heart's electrical system
    \item \textbf{Genetic testing:} for inherited conditions and family screening
    \item \textbf{Pre-procedure assessment:} fitness, medicines, and planning
\end{itemize}

\section*{Management}
\begin{itemize}
    \item \textbf{Insertion procedure:} the generator is placed under the skin and leads advanced into the heart through veins
    \item \textbf{Device programming and testing:} induction of a brief arrhythmia during the procedure to test the shock
    \item \textbf{Aftercare:} wound care, arm restrictions, and simple analgesia
    \item \textbf{Follow-up:} device checks in clinic or remotely, with review of events and battery
    \item \textbf{Optimising heart failure treatment:} medicines to reduce arrhythmia burden
    \item \textbf{Addressing shocks:} adjusting the device and medicines to minimise shocks while staying safe
    \item \textbf{Subcutaneous ICD:} an alternative placed under the skin without leads in the heart, for selected people
    \item \textbf{Psychological support:} counselling for anxiety about shocks and the device
\end{itemize}

\section*{Complications}
\begin{itemize}
    \item Infection of the device or leads, requiring antibiotics and sometimes removal
    \item Bleeding, bruising, and blood collections at the site
    \item Lead displacement or fracture
    \item Pneumothorax from the vein puncture
    \item Inappropriate shocks, delivered for rhythms that are not dangerous
    \item Device-related anxiety and reduced quality of life in some people
    \item Rarely, perforation of the heart or major blood vessels
\end{itemize}

\section*{Prognosis and Outcomes}
ICDs save lives. In people with severe heart failure, implantation reduces the risk of sudden death by around 20 to 30 per cent, and in survivors of cardiac arrest the protection is even greater. Most recipients never receive a life-saving shock, but the device provides reassurance and safety. Shocks occur in a minority, and appropriate device programming and medicines reduce them. With modern devices and follow-up, most people with ICDs live active lives, with restrictions mainly around contact sport and strong magnetic fields.

\section*{Prevention and Self-Care}
\begin{itemize}
    \item Attend all device checks and respond to remote monitoring alerts
    \item Keep the wound clean and report redness, swelling, or discharge
    \item Restrict arm movement on the device side for several weeks
    \item Carry your ICD identification card
    \item Tell all doctors and dentists about the device
    \item Avoid strong magnetic fields and contact sport as advised
    \item Keep mobile phones away from the device
    \item Learn how to respond if the device shocks
    \item Seek psychological support if worried about shocks
    \item Discuss end-of-life care and device deactivation when appropriate
\end{itemize}

\section*{Sources and Further Reading}
The following reputable sources provide further information for healthcare students and patients seeking reliable, current advice about ICD implantation.
\begin{itemize}
    \item Heart Foundation Australia. \textit{Implantable defibrillators.} https://www.heartfoundation.org.au/
    \item Healthdirect Australia. \textit{Implantable cardioverter defibrillator.} https://www.healthdirect.gov.au/
    \item Cardiac Society of Australia and New Zealand. \textit{Device therapy information.} https://www.csanz.edu.au/
    \item Arrhythmia Alliance Australia. \textit{Heart rhythm information.} https://www.heartrhythmalliance.org/
    \item NHS. \textit{Implantable cardioverter defibrillators.} https://www.nhs.uk/conditions/implantable-cardioverter-defibrillators/
    \item Mayo Clinic. \textit{Implantable cardioverter defibrillators.} https://www.mayoclinic.org/
\end{itemize}
